\documentclass[10pt,a4paper]{amsart}
\usepackage[margin=19mm]{geometry}
\usepackage{amsmath,amssymb,mathtools}
\usepackage[scr=rsfs,cal=euler]{mathalfa}
\usepackage{xfrac}
\usepackage[unicode,hidelinks,bookmarks=false]{hyperref}
\newcommand{\ie}{i.e.\ }
\newcommand{\cf}{cf.\ }
\newcommand{\resp}{respectively\ }
\newcommand{\Subsection}{\subsection}
\providecommand{\nobreakdash}{\nobreak\hbox{-}}
\setlength{\emergencystretch}{2em}
\multlinegap0pt
\usepackage{amssymb}
\usepackage{mathtools}
\usepackage[shortlabels]{enumitem}
\usepackage{xcolor}
\newtheorem{proposition}{Proposition}
\newcommand{\Norm}{\operatorname{N}}
\newcommand{\OK}{\mathcal{O}_K}
\newcommand{\Q}{\mathbb{Q}}
\newcommand{\R}{\mathbb{R}}
\newcommand{\dist}{\operatorname{dist}}
\title{The sharp exponent for the minimal distance problem}
\author{Cosmin Pohoata}
\date{Source: arXiv:2607.20422v2 (CC BY 4.0)}
\begin{document}
\maketitle

\noindent\emph{Source and licence.} The sharp exponent for the minimal distance problem. Version arXiv:2607.20422v2 (CC BY 4.0), distributed by arXiv under the Creative Commons Attribution 4.0 International licence, http://creativecommons.org/licenses/by/4.0/, as recorded in the arXiv licence metadata for this exact version and shown on the abstract page https://arxiv.org/abs/2607.20422v2. Full licence notice: \texttt{licenses/arxiv-2607.20422-CC-BY-4.0.txt}.\par\smallskip

\noindent\textit{Editorial scope.} Counted unit: the article's proposition prop:nf-lift (source0.tex lines 616--644). The article's complete original proof of it is delivered with it (source0.tex lines 646--734, one \texttt{proof} environment of 89 lines). No mathematics is written, repaired or re-derived: the statement and the proof are the article's own lines, in the article's own notation, and no retained line is substituted or re-flowed.  Retained uncounted as the source-local context the proof needs: lines 296--299.  Nothing was omitted from any retained span, so the package carries no omission record.  No retained line contains a citation, so the wrapper carries no bibliography and no bibliography entry.  No retained line contains a figure environment, an image-inclusion command or drawing code, so no image is delivered and the package declares no asset dependency.  Editorial formatting only, in addition: an AMS \texttt{amsart} preamble that restates the article's own declarations and macros used by the retained lines, namely the environments \texttt{proposition} on separate counters, together with the standard packages those lines need.  The retained environments print in this document as Proposition 1 in order of appearance, whereas the article prints the same items with the numbering of its own full document.  The complete native source of the article as distributed in the arXiv e-print for this exact version is bundled unchanged as \texttt{sources/205810/source0.tex}.
\begin{equation}
\label{eq:D-definition}
        D(p,p'):=x-x'+y(y'-y).
\end{equation}

\begin{proposition}
\label{prop:nf-lift}
Let $A\subset2\OK\cap B_K(R)$ be square-difference-free, and let
$Y\subset2\OK\cap B_K(M)$, where $R,M\ge1$. Define
\[
        \mathcal P(A,Y)
        :=
        \left\{
        \left(\frac{a+y^2}{2},y\right):
        a\in A,\ y\in Y
        \right\}
        \subset\OK^2,
\]
and put
\begin{equation}
\label{eq:N-H}
        N:=\frac{R+M^2}{2},
        \qquad
        H:=R+2M^2.
\end{equation}
Then one can associate to each $p\in\mathcal P(A,Y)$ a point
$p^*\in[0,1]^2$ and a real line $\ell_p^*$ through $p^*$ such that,
for distinct $p,p'\in\mathcal P(A,Y)$,
\[
        \dist(p'^*,\ell_p^*)
        \ge
        \frac{H^{-(d-1)}}{2\sqrt{N^2+M^4}}.
\]
\end{proposition}

\begin{proof}
The parametrization defining $\mathcal P(A,Y)$ is injective, and hence
$|\mathcal P(A,Y)|=|A||Y|$. Let
\[
        p=(x,y)
        =\left(\frac{a+y^2}{2},y\right)
        \in\mathcal P(A,Y).
\]
Since $a,y\in2\OK$, we have $x\in\OK$. For every real embedding $\rho:K\hookrightarrow\R$, we therefore have
\[
        |\rho(x)|
        \le\frac{|\rho(a)|+|\rho(y)|^2}{2}
        \le\frac{R+M^2}{2}
        =N.
\]
So $(\sigma(x),\sigma(y))$ lies in
$[-N,N]\times[-M,M]$.

Let $L_p^\sigma$ be the real line through
$(\sigma(x),\sigma(y))$ with direction $(\sigma(y),1)$, and define
\begin{equation}
\label{eq:Phi-MN}
        \Phi_{M,N}(X,Y)
        :=\left(\frac{Y+M}{2M},\frac{X+N}{2N}\right).
\end{equation}
This is the symmetric-box version of the previous anisotropic scaling
$(x,y)\mapsto(y/M,x/N)$.  Set
\begin{equation}
\label{eq:scaled-pairs}
        p^*:=\Phi_{M,N}(\sigma(x),\sigma(y)),
        \qquad
        \ell_p^*:=\Phi_{M,N}(L_p^\sigma).
\end{equation}
Then $p^*\in[0,1]^2$ and $p^*\in\ell_p^*$.

Take $p=(x,y)$ and $p'=(x',y')$ in $\mathcal P(A,Y)$, and write
$a=Q(p)$ and $a'=Q(p')$.  The line-equation value from
\eqref{eq:D-definition} is
\begin{equation}
\label{eq:algebraic-D}
\begin{aligned}
        D(p,p') =x-x'+y(y'-y)=\frac{a-a'-(y'-y)^2}{2}\in\OK.
\end{aligned}
\end{equation}
The inclusion in $\OK$ follows from $x,x',y,y'\in\OK$.  If
$D(p,p')=0$, then $a-a'=(y'-y)^2$. Square-difference-freeness gives $a=a'$, hence $y'=y$, and then
$x'=x$.  Therefore
\begin{equation}
\label{eq:D-nonzero}
        p\ne p'
        \quad\Longrightarrow\quad
        D(p,p')\ne0.
\end{equation}

For every real embedding $\rho$, the second expression in
\eqref{eq:algebraic-D} gives
\[
        |\rho(D(p,p'))|
        \le
        \frac{|\rho(a-a')|+|\rho(y'-y)|^2}{2}
        \le R+2M^2
        =H.
\]
When $p\ne p'$, the algebraic integer $D(p,p')$ is nonzero.  Its norm is
therefore a nonzero rational integer, and
\[
        1
        \le|\Norm_{K/\Q}(D(p,p'))|
        =|\sigma(D(p,p'))|
         \prod_{\rho\ne\sigma}|\rho(D(p,p'))|
        \le|\sigma(D(p,p'))|H^{d-1}.
\]
Consequently,
\begin{equation}
\label{eq:D-lower}
        |\sigma(D(p,p'))|\ge H^{-(d-1)}.
\end{equation}

Under $\Phi_{M,N}$, it is easy to check that the direction vector for the line $\ell_p^*$ is given by $v_p=\left(\frac1{2M},\frac{\sigma(y)}{2N}\right)$,
whereas $p'^*-p^* =\left( \frac{\sigma(y'-y)}{2M}, \frac{\sigma(x'-x)}{2N}\right)$. The determinant formula for point-to-line distance then yields the exact identity
\begin{equation}
\label{eq:nf-distance-exact}
\begin{aligned}
        \dist(p'^*,\ell_p^*)=\frac{|\det(p'^*-p^*,v_p)|}{\|v_p\|_2}=\frac{|\sigma(D(p,p'))|}
                {2\sqrt{N^2+M^2\sigma(y)^2}}.
\end{aligned}
\end{equation}
Combining \eqref{eq:D-lower} with $|\sigma(y)|\le M$ proves Proposition \ref{prop:nf-lift}.
\end{proof}

\end{document}
